\section{Códigos do servidor}\label{ch:codigos-servidor}

\subsection{Código do \textit{schema} do endpoint de
status}\label{schema:codigo-schema-status}
\begin{lstlisting}[language=python]
from pydantic import BaseModel

class CpuFreq(BaseModel):
    current: float
    min: float
    max: float

class Memory(BaseModel):
    total: int
    available: int
    percent: float
    used: int
    free: int
    active: int
    inactive: int
    buffers: int
    cached: int
    shared: int
    slab: int

class StatusPerCore(BaseModel):
    cpu_frequency: CpuFreq
    cpu_percent: float

class StatusSchema(BaseModel):
    number_of_cores: int
    cpu_frequency: CpuFreq
    cpu_percent: float
    status_per_core: list[StatusPerCore]
    memory: Memory
\end{lstlisting}

\subsection{Código que carrega variáveis de ambiente
automaticamente}\label{schema:codigo-load-env}
\begin{lstlisting}[language=python]
from functools import lru_cache
from pydantic_settings import BaseSettings, SettingsConfigDict

class Settings(BaseSettings):
    model_config = SettingsConfigDict(env_file='.env', env_file_encoding='utf-8', extra='ignore')

    REDIS_HOST: str = 'localhost'
    REDIS_PORT: int = 6379

@lru_cache
def get_settings():
    return Settings()
\end{lstlisting}

\subsection{Código da função de hash de \acsp{ID}}\label{hash:codigo-hash-id}
\begin{lstlisting}[language=python]
import hashlib

def hash_id(user_id: str) -> str:
    return hashlib.sha256(user_id.encode('utf-8')).hexdigest()
\end{lstlisting}
